\documentclass[11pt]{article}

\usepackage[margin=1in]{geometry}
\usepackage[T1]{fontenc}
\usepackage[utf8]{inputenc}
\usepackage{lmodern}
\usepackage{enumitem}
\usepackage{hyperref}

\hypersetup{
  colorlinks=true,
  linkcolor=blue!50!black,
  urlcolor=blue!50!black,
  citecolor=blue!50!black,
  pdftitle={Discovery Pattern Routing for p-vs-np}
}

\title{Discovery Pattern Routing for \texttt{p-vs-np}}
\author{Neuro-symbolic AI Group\\Idiap Research Institute}
\date{March 18, 2026}

\begin{document}
\maketitle

\begin{abstract}
This report turns blocked-frontier prompt packs into an executable discovery-routing layer. It spawns dedicated discovery layers per breakthrough family and per cross-cutting epistemic pattern, then orders them to avoid collapsing into one depth-first search line.
\end{abstract}

\section{Summary}
\begin{itemize}[leftmargin=*,itemsep=0.2em]
\item Discovery layers are now spawned from prompt packs rather than stored only as passive prompt artifacts.
\item Routing alternates breakthrough-family branches with cross-cutting epistemic-pattern branches to avoid indeterminate depth-first search.
\item Each layer records its own theorem target, stop criteria, and expected deliverables so discovery remains inspectable and finite.
\item Each layer also carries a local-first heuristic for when to search externally for a mathematical framework with similar structural properties.
\end{itemize}

\section{Routing Policy}
\begin{itemize}[leftmargin=*,itemsep=0.2em]
\item Prioritize pressure on the blocked frontier, but enforce diversity across layer kinds and breakthrough families.
\item Use family-specific layers to attack concrete bridge defects and epistemic-pattern layers to change representation, invariants, or decomposition when local search stalls.
\item Escalate or stop when a layer only reproduces its anti-patterns or fails to yield a theorem-shaped claim.
\item Trigger external framework search only after a recorded local failure, and only when the missing move is better modeled as a borrowed foundation than as another local lemma.
\end{itemize}

\section{Dispatch Rounds}
\subsection*{Round 1}
\textbf{Rationale.} Alternate targeted family branches with broader epistemic patterns to avoid collapsing into one depth-first line of search.

\textbf{Layers.} disclayer.p-vs-np.pattern.invariant-discovery, disclayer.p-vs-np.family.abstraction-search, disclayer.p-vs-np.pattern.regime-localization

\textbf{Covered breakthrough families.} abstraction-search, realization, transport, bridge-specificity, barrier-aware transfer, alignment, explicitness, constructive instantiation, semantic bridge
\subsection*{Round 2}
\textbf{Rationale.} Alternate targeted family branches with broader epistemic patterns to avoid collapsing into one depth-first line of search.

\textbf{Layers.} disclayer.p-vs-np.family.realization, disclayer.p-vs-np.pattern.hybridization-of-partial-bridges, disclayer.p-vs-np.family.bridge-specificity

\textbf{Covered breakthrough families.} realization, abstraction-search, transport, bridge-specificity, explicitness, constructive instantiation, semantic bridge
\subsection*{Round 3}
\textbf{Rationale.} Alternate targeted family branches with broader epistemic patterns to avoid collapsing into one depth-first line of search.

\textbf{Layers.} disclayer.p-vs-np.pattern.counterexample-led-concept-formation, disclayer.p-vs-np.family.transport, disclayer.p-vs-np.pattern.route-decomposition

\textbf{Covered breakthrough families.} abstraction-search, realization, transport, barrier-aware transfer, alignment, semantic bridge, bridge-specificity
\subsection*{Round 4}
\textbf{Rationale.} Alternate targeted family branches with broader epistemic patterns to avoid collapsing into one depth-first line of search.

\textbf{Layers.} disclayer.p-vs-np.family.barrier-aware-transfer, disclayer.p-vs-np.pattern.problem-representation-shift, disclayer.p-vs-np.family.alignment

\textbf{Covered breakthrough families.} barrier-aware transfer, abstraction-search, transport, alignment, semantic bridge
\subsection*{Round 5}
\textbf{Rationale.} Alternate targeted family branches with broader epistemic patterns to avoid collapsing into one depth-first line of search.

\textbf{Layers.} disclayer.p-vs-np.family.explicitness, disclayer.p-vs-np.family.constructive-instantiation, disclayer.p-vs-np.family.semantic-bridge

\textbf{Covered breakthrough families.} explicitness, constructive instantiation, semantic bridge
\subsection*{Round 6}
\textbf{Rationale.} Alternate targeted family branches with broader epistemic patterns to avoid collapsing into one depth-first line of search.

\textbf{Layers.} disclayer.p-vs-np.family.categorical-reframing, disclayer.p-vs-np.family.equivalence-coherence-reframing

\textbf{Covered breakthrough families.} realization, transport, bridge-specificity, barrier-aware transfer, alignment, constructive instantiation, semantic bridge, explicitness
\section{disclayer.p-vs-np.pattern.invariant-discovery}
\textbf{Kind.} epistemic-pattern

\textbf{Target.} invariant-discovery

\textbf{Routing priority.} 6.868

\textbf{Dispatch rank.} 1

\textbf{Spawn reason.} Spawn 'invariant-discovery' as a cross-cutting discovery layer because the blocked frontier spans multiple breakthrough families and needs a broader reframing.

\textbf{Rationale.} Several live routes are still searching for the right preserved quantity or witness measure, so invariant discovery remains a first-class discovery move.

\subsection*{System Prompt}
You are searching for the invariant that the current route is really trying to preserve, expose, or measure. Do not begin from the named bridge object. Begin from the quantity, structure, or monotone property that would survive the route and explain the payoff.

\subsection*{User Task}
Identify 2-4 candidate invariants, state how each is computed or witnessed, and say which one would most sharply reorganize the current blocked route.

\subsection*{Expected Deliverables}
\begin{itemize}[leftmargin=*,itemsep=0.2em]
\item Produce one reframing under 'invariant-discovery' that changes the search state rather than restating the blocked lemma.
\item Connect the reframing back to a theorem-shaped objective compatible with 'There exists an admissible MCSP witness family, realized as a family of instances equipped with a monotone witness measure, such that whenever A barely-nontrivial lower bound for an explicit problem Q in a weak model holds for MCSP in nonuniform-circuits, the object organizes the source-side mathematical content while preserving the invariant MCSP complexity profile in nonuniform-circuits and yields major consequences such as NP not in P/poly.'.
\item Record one stop condition showing when the pattern has become superficial rather than discovery-bearing.
\end{itemize}
\subsection*{Stop Criteria}
\begin{itemize}[leftmargin=*,itemsep=0.2em]
\item Stop immediately if the branch collapses into the anti-pattern 'accepting existence without algorithmic access'.
\item Stop if the branch yields no theorem-shaped claim, carrier, transport principle, invariant, or representation change.
\item Stop after one explicit positive candidate and one explicit failure probe have both been recorded.
\item Escalate if 'invariant-discovery' only re-labels the blocked frontier instead of reducing proof debt or changing the representation.
\end{itemize}
\subsection*{External Framework Search Heuristic}
\textbf{Recommended.} True

\textbf{Mode.} targeted-foundation-search

\textbf{Justification.} Search externally for a mathematical framework with similar structural properties only after local routes, InformalMathLib, and explicit negative probes fail to provide the missing carrier, invariant, transfer law, or correspondence theorem.

\textbf{Triggers.} structural-analogy, cross-family-pressure, pattern-level-reframing, bridge-family-foundation-gap

\textbf{Candidate framework regions.} potential / energy method frameworks, monotone measure frameworks, certificate-complexity and witness-measure theories, proof complexity, pseudorandomness / derandomization

\subsection*{Local Prerequisites Before Browsing}
\begin{itemize}[leftmargin=*,itemsep=0.2em]
\item Record the failed local theorem lookup, InformalMathLib lookup, and negative probe that justify leaving local search.
\item State the missing operation concretely: carrier realization, transport, alignment, explicitness, characterization, or invariant discovery.
\item Reject at least one purely topical analogy before searching for an external framework.
\end{itemize}
\subsection*{Query Templates}
\begin{itemize}[leftmargin=*,itemsep=0.2em]
\item mathematical framework for Invariant Discovery Prompt with preserved invariant and bridge theorem
\item Invariant Discovery Prompt correspondence theorem transport invariant
\item Invariant Discovery Prompt framework Several live routes are still searching for the right preserved quantity or witness measure, so invariant discovery remains a first-class discovery move. theorem
\item theorem or framework matching: There exists an admissible MCSP witness family, realized as a family of instances equipped with a monotone witness measure, such that whenever A barely-nontrivial lower bound for an explicit problem Q in a weak model holds for MCSP in nonuniform-circuits, the object organizes the source-side mathematical content while preserving the invariant MCSP complexity profile in nonuniform-circuits and yields major consequences such as NP not in P/poly.
\end{itemize}
\subsection*{Required Questions}
\begin{itemize}[leftmargin=*,itemsep=0.2em]
\item What quantity or structure must survive every legitimate route variant?
\item Which candidate invariant is strong enough to force the payoff but weak enough to be realizable?
\item How would failure of the invariant candidate sharpen the obstruction class?
\end{itemize}
\subsection*{Search Modes}
\begin{itemize}[leftmargin=*,itemsep=0.2em]
\item activation theorem search
\item barrier-aware gadget design
\item bridge split
\item bridge-localization refinement
\item canonical example construction
\item certificate extraction
\item certificate map construction
\item counterfactual sharpening
\item cross-model transport principle
\item dual characterization of the target object
\item duality or normal form search
\item effective construction
\item equivalence search between live carrier encodings
\item example-preserving reduction search
\item explicit carrier construction
\item explicit extractor or generator search
\item generalization by faithful embedding
\item hypothesis verification search
\item impossibility theorem search
\item invariant-preserving compilation theorem
\item normal-form mediation
\item object-class correspondence search
\item parameter-preserving reformulation
\item problem-specific to generic theorem bridge
\item reduction regime decomposition
\item representation normal form for formula-size transfer
\item semantic correspondence theorem
\item semantic-to-algorithmic upgrade theorem
\item surrogate problem design with forcing semantics
\item template instantiation
\item two-way implication decomposition
\item typed abstraction search
\item uniformization
\end{itemize}
\section{disclayer.p-vs-np.family.abstraction-search}
\textbf{Kind.} breakthrough-family

\textbf{Target.} abstraction-search

\textbf{Routing priority.} 6.443

\textbf{Dispatch rank.} 2

\textbf{Spawn reason.} Spawn a dedicated 'abstraction-search' discovery layer because obstruction 'mcsp-magnification' carries repeated pressure across the blocked frontier.

\textbf{Rationale.} Find a sharper abstraction that turns the blocked route into a typed, inspectable theorem target.

\subsection*{System Prompt}
You are searching for the minimal conceptual breakthrough that would turn a blocked route into a typed theorem target. Do not restate the blocked lemma. Identify the sharper carrier, invariant, or bridge law that would make the next theorem materially more specific.

\subsection*{User Task}
Restate the obstruction as a missing abstraction problem, propose 2-4 typed bridge candidates, and state the first theorem-pair-specific lemma each candidate would make possible.

\subsection*{Expected Deliverables}
\begin{itemize}[leftmargin=*,itemsep=0.2em]
\item State one theorem-shaped candidate centered on 'abstraction-search' satisfying 'State a route-specific theorem that would resolve the blocked interface and make the next proof obligation materially sharper.'.
\item Make 'A theorem-family-specific carrier, invariant, or bridge law that adds new mathematical content.' explicit enough that an expert could inspect the object or transfer law.
\item Record one decisive negative route explaining what would fail if this discovery layer does not land.
\end{itemize}
\subsection*{Stop Criteria}
\begin{itemize}[leftmargin=*,itemsep=0.2em]
\item Stop immediately if the branch collapses into the anti-pattern 'repeating the same blocked statement without new typed structure'.
\item Stop if the branch yields no theorem-shaped claim, carrier, transport principle, invariant, or representation change.
\item Stop after one explicit positive candidate and one explicit failure probe have both been recorded.
\item Escalate if 'abstraction-search' remains purely schematic after the first explicit object or theorem candidate.
\end{itemize}
\subsection*{External Framework Search Heuristic}
\textbf{Recommended.} True

\textbf{Mode.} targeted-foundation-search

\textbf{Justification.} Search externally for a mathematical framework with similar structural properties only after local routes, InformalMathLib, and explicit negative probes fail to provide the missing carrier, invariant, transfer law, or correspondence theorem.

\textbf{Triggers.} cross-family-pressure, bridge-family-foundation-gap

\textbf{Candidate framework regions.} proof complexity, pseudorandomness / derandomization, coding-theoretic hardness surrogates

\subsection*{Local Prerequisites Before Browsing}
\begin{itemize}[leftmargin=*,itemsep=0.2em]
\item Record the failed local theorem lookup, InformalMathLib lookup, and negative probe that justify leaving local search.
\item State the missing operation concretely: carrier realization, transport, alignment, explicitness, characterization, or invariant discovery.
\item Reject at least one purely topical analogy before searching for an external framework.
\end{itemize}
\subsection*{Query Templates}
\begin{itemize}[leftmargin=*,itemsep=0.2em]
\item mathematical framework for Abstraction Search Prompt with preserved invariant and bridge theorem
\item Abstraction Search Prompt correspondence theorem transport invariant
\item Abstraction Search Prompt framework Find a sharper abstraction that turns the blocked route into a typed, inspectable theorem target. theorem
\item theorem or framework matching: State a route-specific theorem that would resolve the blocked interface and make the next proof obligation materially sharper.
\end{itemize}
\subsection*{Required Questions}
\begin{itemize}[leftmargin=*,itemsep=0.2em]
\item What typed object would add mathematical content beyond the current schematic carrier?
\item What route-specific theorem would become stateable if that object existed?
\item Which candidate clarifies the negative route if it fails?
\end{itemize}
\subsection*{Search Modes}
\begin{itemize}[leftmargin=*,itemsep=0.2em]
\item typed abstraction search
\item bridge split
\item counterfactual sharpening
\end{itemize}
\section{disclayer.p-vs-np.pattern.regime-localization}
\textbf{Kind.} epistemic-pattern

\textbf{Target.} regime-localization

\textbf{Routing priority.} 6.758

\textbf{Dispatch rank.} 3

\textbf{Spawn reason.} Spawn 'regime-localization' as a cross-cutting discovery layer because the blocked frontier spans multiple breakthrough families and needs a broader reframing.

\textbf{Rationale.} The frontier repeatedly crosses model or threshold boundaries, so discovery should search for the smallest regime in which the route first becomes true.

\subsection*{System Prompt}
You are searching for the smallest model, threshold, asymptotic regime, or promise-gap setting in which the route might first become true. Do not ask for the full theorem immediately if a more local regime would already be a major contribution.

\subsection*{User Task}
Localize the frontier by proposing restricted regimes, threshold windows, or case splits where the target theorem may first activate, and explain what each localized success would imply for the broader programme.

\subsection*{Expected Deliverables}
\begin{itemize}[leftmargin=*,itemsep=0.2em]
\item Produce one reframing under 'regime-localization' that changes the search state rather than restating the blocked lemma.
\item Connect the reframing back to a theorem-shaped objective compatible with 'There exists an admissible MCSP witness family, realized as a family of instances equipped with a monotone witness measure, such that whenever A barely-nontrivial lower bound for an explicit problem Q in a weak model holds for MCSP in nonuniform-circuits, the object organizes the source-side mathematical content while preserving the invariant MCSP complexity profile in nonuniform-circuits and yields major consequences such as NP not in P/poly.'.
\item Record one stop condition showing when the pattern has become superficial rather than discovery-bearing.
\end{itemize}
\subsection*{Stop Criteria}
\begin{itemize}[leftmargin=*,itemsep=0.2em]
\item Stop immediately if the branch collapses into the anti-pattern 'adding more obligations without exhibiting a concrete carrier family'.
\item Stop if the branch yields no theorem-shaped claim, carrier, transport principle, invariant, or representation change.
\item Stop after one explicit positive candidate and one explicit failure probe have both been recorded.
\item Escalate if 'regime-localization' only re-labels the blocked frontier instead of reducing proof debt or changing the representation.
\end{itemize}
\subsection*{External Framework Search Heuristic}
\textbf{Recommended.} True

\textbf{Mode.} targeted-foundation-search

\textbf{Justification.} Search externally for a mathematical framework with similar structural properties only after local routes, InformalMathLib, and explicit negative probes fail to provide the missing carrier, invariant, transfer law, or correspondence theorem.

\textbf{Triggers.} structural-analogy, cross-family-pressure, pattern-level-reframing, bridge-family-foundation-gap

\textbf{Candidate framework regions.} threshold and regime-restriction frameworks, restricted-model lower-bound programmes, proof complexity, pseudorandomness / derandomization, coding-theoretic hardness surrogates

\subsection*{Local Prerequisites Before Browsing}
\begin{itemize}[leftmargin=*,itemsep=0.2em]
\item Record the failed local theorem lookup, InformalMathLib lookup, and negative probe that justify leaving local search.
\item State the missing operation concretely: carrier realization, transport, alignment, explicitness, characterization, or invariant discovery.
\item Reject at least one purely topical analogy before searching for an external framework.
\end{itemize}
\subsection*{Query Templates}
\begin{itemize}[leftmargin=*,itemsep=0.2em]
\item mathematical framework for Regime Localization Prompt with preserved invariant and bridge theorem
\item Regime Localization Prompt correspondence theorem transport invariant
\item Regime Localization Prompt framework The frontier repeatedly crosses model or threshold boundaries, so discovery should search for the smallest regime in which the route first becomes true. theorem
\item theorem or framework matching: There exists an admissible MCSP witness family, realized as a family of instances equipped with a monotone witness measure, such that whenever A barely-nontrivial lower bound for an explicit problem Q in a weak model holds for MCSP in nonuniform-circuits, the object organizes the source-side mathematical content while preserving the invariant MCSP complexity profile in nonuniform-circuits and yields major consequences such as NP not in P/poly.
\end{itemize}
\subsection*{Required Questions}
\begin{itemize}[leftmargin=*,itemsep=0.2em]
\item Which regime change most plausibly removes the current obstruction without trivializing the result?
\item What is the first nontrivial threshold at which the route may become true?
\item Would a localized theorem already count as an independently valuable subresult?
\end{itemize}
\subsection*{Search Modes}
\begin{itemize}[leftmargin=*,itemsep=0.2em]
\item activation theorem search
\item barrier-aware gadget design
\item bridge split
\item bridge-localization refinement
\item canonical example construction
\item counterfactual sharpening
\item cross-model transport principle
\item equivalence search between live carrier encodings
\item example-preserving reduction search
\item explicit carrier construction
\item generalization by faithful embedding
\item hypothesis verification search
\item impossibility theorem search
\item invariant-preserving compilation theorem
\item parameter-preserving reformulation
\item problem-specific to generic theorem bridge
\item reduction regime decomposition
\item representation normal form for formula-size transfer
\item surrogate problem design with forcing semantics
\item template instantiation
\item typed abstraction search
\end{itemize}
\section{disclayer.p-vs-np.family.realization}
\textbf{Kind.} breakthrough-family

\textbf{Target.} realization

\textbf{Routing priority.} 6.443

\textbf{Dispatch rank.} 4

\textbf{Spawn reason.} Spawn a dedicated 'realization' discovery layer because obstruction 'carrier-realization' carries repeated pressure across the blocked frontier.

\textbf{Rationale.} Turn a typed but schematic carrier into an explicit mathematical object family with admissibility conditions and a forcing theorem.

\subsection*{System Prompt}
You are converting a typed but schematic carrier into an explicit mathematical object family. Do not rename the carrier. Produce admissibility conditions, a realizability route from the source theorem, and the forcing lemma needed at the target.

\subsection*{User Task}
Construct or sharply constrain an explicit witness family, obstruction family, surrogate promise problem, or quantitative measure that realizes the blocked route from source theorem to target theorem.

\subsection*{Expected Deliverables}
\begin{itemize}[leftmargin=*,itemsep=0.2em]
\item State one theorem-shaped candidate centered on 'realization' satisfying 'If source theorem S holds in regime R, then one can construct an explicit carrier C with admissibility conditions A such that C preserves invariant I and supports a route-specific forcing lemma to target T.'.
\item Make 'An explicit obstruction family, witness family, surrogate promise problem, or quantitative measure that is realized from the source theorem rather than only named by role.' explicit enough that an expert could inspect the object or transfer law.
\item Record one decisive negative route explaining what would fail if this discovery layer does not land.
\end{itemize}
\subsection*{Stop Criteria}
\begin{itemize}[leftmargin=*,itemsep=0.2em]
\item Stop immediately if the branch collapses into the anti-pattern 'renaming the carrier without adding realizability conditions'.
\item Stop if the branch yields no theorem-shaped claim, carrier, transport principle, invariant, or representation change.
\item Stop after one explicit positive candidate and one explicit failure probe have both been recorded.
\item Escalate if 'realization' remains purely schematic after the first explicit object or theorem candidate.
\end{itemize}
\subsection*{External Framework Search Heuristic}
\textbf{Recommended.} True

\textbf{Mode.} targeted-foundation-search

\textbf{Justification.} Search externally for a mathematical framework with similar structural properties only after local routes, InformalMathLib, and explicit negative probes fail to provide the missing carrier, invariant, transfer law, or correspondence theorem.

\textbf{Triggers.} structural-analogy, cross-family-pressure, bridge-family-foundation-gap

\textbf{Candidate framework regions.} proof complexity witness measures, coding theory, pseudorandom objects and hardness amplification

\subsection*{Local Prerequisites Before Browsing}
\begin{itemize}[leftmargin=*,itemsep=0.2em]
\item Record the failed local theorem lookup, InformalMathLib lookup, and negative probe that justify leaving local search.
\item State the missing operation concretely: carrier realization, transport, alignment, explicitness, characterization, or invariant discovery.
\item Reject at least one purely topical analogy before searching for an external framework.
\end{itemize}
\subsection*{Query Templates}
\begin{itemize}[leftmargin=*,itemsep=0.2em]
\item mathematical framework for Carrier Realization Prompt with preserved invariant and bridge theorem
\item Carrier Realization Prompt correspondence theorem transport invariant
\item Carrier Realization Prompt framework Turn a typed but schematic carrier into an explicit mathematical object family with admissibility conditions and a forcing theorem. theorem
\item theorem or framework matching: If source theorem S holds in regime R, then one can construct an explicit carrier C with admissibility conditions A such that C preserves invariant I and supports a route-specific forcing lemma to target T.
\end{itemize}
\subsection*{Required Questions}
\begin{itemize}[leftmargin=*,itemsep=0.2em]
\item What is the concrete carrier set, map, family, or measure?
\item How is it built from the source theorem rather than postulated?
\item What admissibility conditions prevent it from being only a relabeling?
\end{itemize}
\subsection*{Search Modes}
\begin{itemize}[leftmargin=*,itemsep=0.2em]
\item explicit carrier construction
\item equivalence search between live carrier encodings
\item surrogate problem design with forcing semantics
\end{itemize}
\section{disclayer.p-vs-np.pattern.hybridization-of-partial-bridges}
\textbf{Kind.} epistemic-pattern

\textbf{Target.} hybridization-of-partial-bridges

\textbf{Routing priority.} 6.728

\textbf{Dispatch rank.} 5

\textbf{Spawn reason.} Spawn 'hybridization-of-partial-bridges' as a cross-cutting discovery layer because the blocked frontier spans multiple breakthrough families and needs a broader reframing.

\textbf{Rationale.} Some blocked routes sit at the intersection of multiple breakthrough families, which is evidence that one bridge object may be insufficient and a composed route should be explored.

\subsection*{System Prompt}
Do not assume one bridge object must do all the work. Search for a hybrid route in which two partial bridges compose, with each bridge carrying a distinct inferential burden.

\subsection*{User Task}
Identify complementary partial bridges, specify their composition order, and state the intermediate theorem that would let them work together.

\subsection*{Expected Deliverables}
\begin{itemize}[leftmargin=*,itemsep=0.2em]
\item Produce one reframing under 'hybridization-of-partial-bridges' that changes the search state rather than restating the blocked lemma.
\item Connect the reframing back to a theorem-shaped objective compatible with 'There exists an admissible MCSP witness family, realized as a family of instances equipped with a monotone witness measure, such that whenever A barely-nontrivial lower bound for an explicit problem Q in a weak model holds for MCSP in nonuniform-circuits, the object organizes the source-side mathematical content while preserving the invariant MCSP complexity profile in nonuniform-circuits and yields major consequences such as NP not in P/poly.'.
\item Record one stop condition showing when the pattern has become superficial rather than discovery-bearing.
\end{itemize}
\subsection*{Stop Criteria}
\begin{itemize}[leftmargin=*,itemsep=0.2em]
\item Stop immediately if the branch collapses into the anti-pattern 'accepting existence without algorithmic access'.
\item Stop if the branch yields no theorem-shaped claim, carrier, transport principle, invariant, or representation change.
\item Stop after one explicit positive candidate and one explicit failure probe have both been recorded.
\item Escalate if 'hybridization-of-partial-bridges' only re-labels the blocked frontier instead of reducing proof debt or changing the representation.
\end{itemize}
\subsection*{External Framework Search Heuristic}
\textbf{Recommended.} True

\textbf{Mode.} targeted-foundation-search

\textbf{Justification.} Search externally for a mathematical framework with similar structural properties only after local routes, InformalMathLib, and explicit negative probes fail to provide the missing carrier, invariant, transfer law, or correspondence theorem.

\textbf{Triggers.} structural-analogy, cross-family-pressure, pattern-level-reframing, bridge-family-foundation-gap

\textbf{Candidate framework regions.} composed transfer frameworks, hybrid argument methods, multi-interface reduction theories, proof complexity, pseudorandomness / derandomization

\subsection*{Local Prerequisites Before Browsing}
\begin{itemize}[leftmargin=*,itemsep=0.2em]
\item Record the failed local theorem lookup, InformalMathLib lookup, and negative probe that justify leaving local search.
\item State the missing operation concretely: carrier realization, transport, alignment, explicitness, characterization, or invariant discovery.
\item Reject at least one purely topical analogy before searching for an external framework.
\end{itemize}
\subsection*{Query Templates}
\begin{itemize}[leftmargin=*,itemsep=0.2em]
\item mathematical framework for Hybrid Bridge Synthesis Prompt with preserved invariant and bridge theorem
\item Hybrid Bridge Synthesis Prompt correspondence theorem transport invariant
\item Hybrid Bridge Synthesis Prompt framework Some blocked routes sit at the intersection of multiple breakthrough families, which is evidence that one bridge object may be insufficient and a composed route should be explored. theorem
\item theorem or framework matching: There exists an admissible MCSP witness family, realized as a family of instances equipped with a monotone witness measure, such that whenever A barely-nontrivial lower bound for an explicit problem Q in a weak model holds for MCSP in nonuniform-circuits, the object organizes the source-side mathematical content while preserving the invariant MCSP complexity profile in nonuniform-circuits and yields major consequences such as NP not in P/poly.
\end{itemize}
\subsection*{Required Questions}
\begin{itemize}[leftmargin=*,itemsep=0.2em]
\item Which blocked routes solve different parts of the same theorem pair?
\item What intermediate object or theorem would allow the partial bridges to compose?
\item What failure would show that the hybrid route is illegitimate rather than merely incomplete?
\end{itemize}
\subsection*{Search Modes}
\begin{itemize}[leftmargin=*,itemsep=0.2em]
\item activation theorem search
\item bridge split
\item bridge-localization refinement
\item canonical example construction
\item certificate extraction
\item certificate map construction
\item counterfactual sharpening
\item cross-model transport principle
\item duality or normal form search
\item effective construction
\item equivalence search between live carrier encodings
\item example-preserving reduction search
\item explicit carrier construction
\item explicit extractor or generator search
\item hypothesis verification search
\item invariant-preserving compilation theorem
\item representation normal form for formula-size transfer
\item semantic correspondence theorem
\item semantic-to-algorithmic upgrade theorem
\item surrogate problem design with forcing semantics
\item template instantiation
\item two-way implication decomposition
\item typed abstraction search
\item uniformization
\end{itemize}
\section{disclayer.p-vs-np.family.bridge-specificity}
\textbf{Kind.} breakthrough-family

\textbf{Target.} bridge-specificity

\textbf{Routing priority.} 5.669

\textbf{Dispatch rank.} 6

\textbf{Spawn reason.} Spawn a dedicated 'bridge-specificity' discovery layer because obstruction 'target-grounding' carries repeated pressure across the blocked frontier.

\textbf{Rationale.} Stop treating the target theorem as the open problem and relocate search to the bridge that would activate it.

\subsection*{System Prompt}
The target theorem is already known. Your task is to identify the local activation theorem that would verify its hypotheses from current artifacts. Do not rediscover the target theorem.

\subsection*{User Task}
State the activation theorem under search, name the exact target hypothesis that is still unverified, and identify the local bridge object or transfer law needed to verify it.

\subsection*{Expected Deliverables}
\begin{itemize}[leftmargin=*,itemsep=0.2em]
\item State one theorem-shaped candidate centered on 'bridge-specificity' satisfying 'Target theorem T is fixed. The real theorem under search is the local activation theorem A that verifies T's hypotheses from current artifacts.'.
\item Make 'No new target theorem; discovery should instead seek the carrier, transfer law, or alignment theorem that makes the target theorem applicable.' explicit enough that an expert could inspect the object or transfer law.
\item Record one decisive negative route explaining what would fail if this discovery layer does not land.
\end{itemize}
\subsection*{Stop Criteria}
\begin{itemize}[leftmargin=*,itemsep=0.2em]
\item Stop immediately if the branch collapses into the anti-pattern 'rediscovering the target theorem'.
\item Stop if the branch yields no theorem-shaped claim, carrier, transport principle, invariant, or representation change.
\item Stop after one explicit positive candidate and one explicit failure probe have both been recorded.
\item Escalate if 'bridge-specificity' remains purely schematic after the first explicit object or theorem candidate.
\end{itemize}
\subsection*{External Framework Search Heuristic}
\textbf{Recommended.} True

\textbf{Mode.} targeted-foundation-search

\textbf{Justification.} Search externally for a mathematical framework with similar structural properties only after local routes, InformalMathLib, and explicit negative probes fail to provide the missing carrier, invariant, transfer law, or correspondence theorem.

\textbf{Triggers.} cross-family-pressure

\textbf{Candidate framework regions.} activation-theorem frameworks, hypothesis-verification templates from neighboring areas

\subsection*{Local Prerequisites Before Browsing}
\begin{itemize}[leftmargin=*,itemsep=0.2em]
\item Record the failed local theorem lookup, InformalMathLib lookup, and negative probe that justify leaving local search.
\item State the missing operation concretely: carrier realization, transport, alignment, explicitness, characterization, or invariant discovery.
\item Reject at least one purely topical analogy before searching for an external framework.
\end{itemize}
\subsection*{Query Templates}
\begin{itemize}[leftmargin=*,itemsep=0.2em]
\item mathematical framework for Activation Theorem Prompt with preserved invariant and bridge theorem
\item Activation Theorem Prompt correspondence theorem transport invariant
\item Activation Theorem Prompt framework Stop treating the target theorem as the open problem and relocate search to the bridge that would activate it. theorem
\item theorem or framework matching: Target theorem T is fixed. The real theorem under search is the local activation theorem A that verifies T's hypotheses from current artifacts.
\end{itemize}
\subsection*{Required Questions}
\begin{itemize}[leftmargin=*,itemsep=0.2em]
\item Which target-theorem hypothesis is still open locally?
\item What bridge-local theorem would certify that hypothesis?
\item What novelty remains once the target theorem itself is removed from search?
\end{itemize}
\subsection*{Search Modes}
\begin{itemize}[leftmargin=*,itemsep=0.2em]
\item activation theorem search
\item hypothesis verification search
\item bridge-localization refinement
\end{itemize}
\section{disclayer.p-vs-np.pattern.counterexample-led-concept-formation}
\textbf{Kind.} epistemic-pattern

\textbf{Target.} counterexample-led-concept-formation

\textbf{Routing priority.} 6.658

\textbf{Dispatch rank.} 7

\textbf{Spawn reason.} Spawn 'counterexample-led-concept-formation' as a cross-cutting discovery layer because the blocked frontier spans multiple breakthrough families and needs a broader reframing.

\textbf{Rationale.} The frontier is heavily blocked, so near-counterexamples and barrier analogs should be used to induce better concepts rather than only to reject claims.

\subsection*{System Prompt}
Use failure as a discovery tool. Search for near-counterexamples, barrier analogs, and route-breaking examples, then ask what concept would separate the surviving cases from the failures.

\subsection*{User Task}
Identify the strongest nearby failure patterns, describe what they invalidate, and derive a refined concept or admissibility condition from those failures.

\subsection*{Expected Deliverables}
\begin{itemize}[leftmargin=*,itemsep=0.2em]
\item Produce one reframing under 'counterexample-led-concept-formation' that changes the search state rather than restating the blocked lemma.
\item Connect the reframing back to a theorem-shaped objective compatible with 'There exists an admissible MCSP witness family, realized as a family of instances equipped with a monotone witness measure, such that whenever A barely-nontrivial lower bound for an explicit problem Q in a weak model holds for MCSP in nonuniform-circuits, the object organizes the source-side mathematical content while preserving the invariant MCSP complexity profile in nonuniform-circuits and yields major consequences such as NP not in P/poly.'.
\item Record one stop condition showing when the pattern has become superficial rather than discovery-bearing.
\end{itemize}
\subsection*{Stop Criteria}
\begin{itemize}[leftmargin=*,itemsep=0.2em]
\item Stop immediately if the branch collapses into the anti-pattern 'adding a stronger carrier without specifying the target object class map'.
\item Stop if the branch yields no theorem-shaped claim, carrier, transport principle, invariant, or representation change.
\item Stop after one explicit positive candidate and one explicit failure probe have both been recorded.
\item Escalate if 'counterexample-led-concept-formation' only re-labels the blocked frontier instead of reducing proof debt or changing the representation.
\end{itemize}
\subsection*{External Framework Search Heuristic}
\textbf{Recommended.} True

\textbf{Mode.} targeted-foundation-search

\textbf{Justification.} Search externally for a mathematical framework with similar structural properties only after local routes, InformalMathLib, and explicit negative probes fail to provide the missing carrier, invariant, transfer law, or correspondence theorem.

\textbf{Triggers.} structural-analogy, cross-family-pressure, bridge-family-foundation-gap

\textbf{Candidate framework regions.} barrier and impossibility frameworks, minimal obstruction and forbidden-configuration theories, proof complexity, pseudorandomness / derandomization, coding-theoretic hardness surrogates

\subsection*{Local Prerequisites Before Browsing}
\begin{itemize}[leftmargin=*,itemsep=0.2em]
\item Record the failed local theorem lookup, InformalMathLib lookup, and negative probe that justify leaving local search.
\item State the missing operation concretely: carrier realization, transport, alignment, explicitness, characterization, or invariant discovery.
\item Reject at least one purely topical analogy before searching for an external framework.
\end{itemize}
\subsection*{Query Templates}
\begin{itemize}[leftmargin=*,itemsep=0.2em]
\item mathematical framework for Counterexample-Led Concept Formation Prompt with preserved invariant and bridge theorem
\item Counterexample-Led Concept Formation Prompt correspondence theorem transport invariant
\item Counterexample-Led Concept Formation Prompt framework The frontier is heavily blocked, so near-counterexamples and barrier analogs should be used to induce better concepts rather than only to reject claims. theorem
\item theorem or framework matching: There exists an admissible MCSP witness family, realized as a family of instances equipped with a monotone witness measure, such that whenever A barely-nontrivial lower bound for an explicit problem Q in a weak model holds for MCSP in nonuniform-circuits, the object organizes the source-side mathematical content while preserving the invariant MCSP complexity profile in nonuniform-circuits and yields major consequences such as NP not in P/poly.
\end{itemize}
\subsection*{Required Questions}
\begin{itemize}[leftmargin=*,itemsep=0.2em]
\item What is the nearest counterexample or barrier analog to the current route?
\item Which current concept fails to distinguish the positive and negative cases?
\item What new concept becomes visible only after inspecting the failure pattern?
\end{itemize}
\subsection*{Search Modes}
\begin{itemize}[leftmargin=*,itemsep=0.2em]
\item barrier-aware gadget design
\item bridge split
\item counterfactual sharpening
\item cross-model transport principle
\item dual characterization of the target object
\item duality or normal form search
\item equivalence search between live carrier encodings
\item explicit carrier construction
\item generalization by faithful embedding
\item impossibility theorem search
\item invariant-preserving compilation theorem
\item normal-form mediation
\item object-class correspondence search
\item parameter-preserving reformulation
\item problem-specific to generic theorem bridge
\item reduction regime decomposition
\item representation normal form for formula-size transfer
\item semantic correspondence theorem
\item surrogate problem design with forcing semantics
\item two-way implication decomposition
\item typed abstraction search
\end{itemize}
\section{disclayer.p-vs-np.family.transport}
\textbf{Kind.} breakthrough-family

\textbf{Target.} transport

\textbf{Routing priority.} 5.610

\textbf{Dispatch rank.} 8

\textbf{Spawn reason.} Spawn a dedicated 'transport' discovery layer because obstruction 'formula-size-transport' carries repeated pressure across the blocked frontier.

\textbf{Rationale.} Find a theorem that preserves the relevant invariant across a cross-model move into formula-size.

\subsection*{System Prompt}
You are searching for an invariant-preserving transport theorem across a model shift. Do not strengthen the source theorem unless it changes the transport mechanism. Focus on the cross-model map and the preserved quantitative parameter.

\subsection*{User Task}
State the strongest plausible transport law, compilation theorem, or normal form that could move the current carrier into the target model without losing the invariant needed for the payoff.

\subsection*{Expected Deliverables}
\begin{itemize}[leftmargin=*,itemsep=0.2em]
\item State one theorem-shaped candidate centered on 'transport' satisfying 'For every admissible carrier C in the nonuniform-circuit regime, there is a transformed carrier C' in the formula-size regime preserving invariant I and retaining enough quantitative strength to imply target T.'.
\item Make 'A transport law, compilation theorem, or complexity measure that is stable under the source-to-target model shift.' explicit enough that an expert could inspect the object or transfer law.
\item Record one decisive negative route explaining what would fail if this discovery layer does not land.
\end{itemize}
\subsection*{Stop Criteria}
\begin{itemize}[leftmargin=*,itemsep=0.2em]
\item Stop immediately if the branch collapses into the anti-pattern 'assuming the target theorem already provides the transport'.
\item Stop if the branch yields no theorem-shaped claim, carrier, transport principle, invariant, or representation change.
\item Stop after one explicit positive candidate and one explicit failure probe have both been recorded.
\item Escalate if 'transport' remains purely schematic after the first explicit object or theorem candidate.
\end{itemize}
\subsection*{External Framework Search Heuristic}
\textbf{Recommended.} True

\textbf{Mode.} targeted-foundation-search

\textbf{Justification.} Search externally for a mathematical framework with similar structural properties only after local routes, InformalMathLib, and explicit negative probes fail to provide the missing carrier, invariant, transfer law, or correspondence theorem.

\textbf{Triggers.} structural-analogy, cross-family-pressure, bridge-family-foundation-gap

\textbf{Candidate framework regions.} formula lower-bound transfer frameworks, communication complexity model-transfer methods, representation/normal-form theorems in circuit complexity

\subsection*{Local Prerequisites Before Browsing}
\begin{itemize}[leftmargin=*,itemsep=0.2em]
\item Record the failed local theorem lookup, InformalMathLib lookup, and negative probe that justify leaving local search.
\item State the missing operation concretely: carrier realization, transport, alignment, explicitness, characterization, or invariant discovery.
\item Reject at least one purely topical analogy before searching for an external framework.
\end{itemize}
\subsection*{Query Templates}
\begin{itemize}[leftmargin=*,itemsep=0.2em]
\item mathematical framework for Transport Principle Prompt with preserved invariant and bridge theorem
\item Transport Principle Prompt correspondence theorem transport invariant
\item Transport Principle Prompt framework Find a theorem that preserves the relevant invariant across a cross-model move into formula-size. theorem
\item theorem or framework matching: For every admissible carrier C in the nonuniform-circuit regime, there is a transformed carrier C' in the formula-size regime preserving invariant I and retaining enough quantitative strength to imply target T.
\end{itemize}
\subsection*{Required Questions}
\begin{itemize}[leftmargin=*,itemsep=0.2em]
\item What exact quantity must be preserved across the model shift?
\item What transformed carrier should live in the target model?
\item What is the first obstruction showing why naive transport fails?
\end{itemize}
\subsection*{Search Modes}
\begin{itemize}[leftmargin=*,itemsep=0.2em]
\item cross-model transport principle
\item invariant-preserving compilation theorem
\item representation normal form for formula-size transfer
\end{itemize}
\section{disclayer.p-vs-np.pattern.route-decomposition}
\textbf{Kind.} epistemic-pattern

\textbf{Target.} route-decomposition

\textbf{Routing priority.} 6.578

\textbf{Dispatch rank.} 9

\textbf{Spawn reason.} Spawn 'route-decomposition' as a cross-cutting discovery layer because the blocked frontier spans multiple breakthrough families and needs a broader reframing.

\textbf{Rationale.} High-pressure obstruction classes are carrying multiple inferential roles at once, so route decomposition is needed before more local search.

\subsection*{System Prompt}
The current route may be carrying too many inferential roles at once. Decompose it into smaller interfaces, each with its own object, invariant, and theorem target.

\subsection*{User Task}
Split the blocked route into a sequence of subinterfaces, identify the minimal theorem at each interface, and state which subresult would already be valuable if the full route remained open.

\subsection*{Expected Deliverables}
\begin{itemize}[leftmargin=*,itemsep=0.2em]
\item Produce one reframing under 'route-decomposition' that changes the search state rather than restating the blocked lemma.
\item Connect the reframing back to a theorem-shaped objective compatible with 'There exists an admissible MCSP witness family, realized as a family of instances equipped with a monotone witness measure, such that whenever A barely-nontrivial lower bound for an explicit problem Q in a weak model holds for MCSP in nonuniform-circuits, the object organizes the source-side mathematical content while preserving the invariant MCSP complexity profile in nonuniform-circuits and yields major consequences such as NP not in P/poly.'.
\item Record one stop condition showing when the pattern has become superficial rather than discovery-bearing.
\end{itemize}
\subsection*{Stop Criteria}
\begin{itemize}[leftmargin=*,itemsep=0.2em]
\item Stop immediately if the branch collapses into the anti-pattern 'adding more obligations without exhibiting a concrete carrier family'.
\item Stop if the branch yields no theorem-shaped claim, carrier, transport principle, invariant, or representation change.
\item Stop after one explicit positive candidate and one explicit failure probe have both been recorded.
\item Escalate if 'route-decomposition' only re-labels the blocked frontier instead of reducing proof debt or changing the representation.
\end{itemize}
\subsection*{External Framework Search Heuristic}
\textbf{Recommended.} True

\textbf{Mode.} targeted-foundation-search

\textbf{Justification.} Search externally for a mathematical framework with similar structural properties only after local routes, InformalMathLib, and explicit negative probes fail to provide the missing carrier, invariant, transfer law, or correspondence theorem.

\textbf{Triggers.} structural-analogy, cross-family-pressure, pattern-level-reframing, bridge-family-foundation-gap

\textbf{Candidate framework regions.} interface decomposition frameworks, multi-stage lifting and factorization methods, proof complexity, pseudorandomness / derandomization, coding-theoretic hardness surrogates

\subsection*{Local Prerequisites Before Browsing}
\begin{itemize}[leftmargin=*,itemsep=0.2em]
\item Record the failed local theorem lookup, InformalMathLib lookup, and negative probe that justify leaving local search.
\item State the missing operation concretely: carrier realization, transport, alignment, explicitness, characterization, or invariant discovery.
\item Reject at least one purely topical analogy before searching for an external framework.
\end{itemize}
\subsection*{Query Templates}
\begin{itemize}[leftmargin=*,itemsep=0.2em]
\item mathematical framework for Route Decomposition Prompt with preserved invariant and bridge theorem
\item Route Decomposition Prompt correspondence theorem transport invariant
\item Route Decomposition Prompt framework High-pressure obstruction classes are carrying multiple inferential roles at once, so route decomposition is needed before more local search. theorem
\item theorem or framework matching: There exists an admissible MCSP witness family, realized as a family of instances equipped with a monotone witness measure, such that whenever A barely-nontrivial lower bound for an explicit problem Q in a weak model holds for MCSP in nonuniform-circuits, the object organizes the source-side mathematical content while preserving the invariant MCSP complexity profile in nonuniform-circuits and yields major consequences such as NP not in P/poly.
\end{itemize}
\subsection*{Required Questions}
\begin{itemize}[leftmargin=*,itemsep=0.2em]
\item Which inferential roles are currently being conflated in one bridge object?
\item What is the smallest theorem that isolates each role?
\item Which decomposed segment is most likely to produce a standalone contribution?
\end{itemize}
\subsection*{Search Modes}
\begin{itemize}[leftmargin=*,itemsep=0.2em]
\item activation theorem search
\item bridge split
\item bridge-localization refinement
\item counterfactual sharpening
\item cross-model transport principle
\item equivalence search between live carrier encodings
\item explicit carrier construction
\item hypothesis verification search
\item invariant-preserving compilation theorem
\item representation normal form for formula-size transfer
\item surrogate problem design with forcing semantics
\item typed abstraction search
\end{itemize}
\section{disclayer.p-vs-np.family.barrier-aware-transfer}
\textbf{Kind.} breakthrough-family

\textbf{Target.} barrier-aware transfer

\textbf{Routing priority.} 2.309

\textbf{Dispatch rank.} 10

\textbf{Spawn reason.} Spawn a dedicated 'barrier-aware transfer' discovery layer because obstruction 'local-reduction-barrier' carries repeated pressure across the blocked frontier.

\textbf{Rationale.} Design a transfer mechanism that survives known local-reduction barriers, or prove that no such mechanism exists under explicit assumptions.

\subsection*{System Prompt}
You are searching under an explicit barrier theorem. Either construct a transfer that survives the barrier's constraints or sharpen the impossibility theorem. Do not use generic hardness language that ignores the reduction regime.

\subsection*{User Task}
Describe a barrier-aware gadget family or impossibility result for the blocked transfer, with locality bounds and preserved witness structure made explicit.

\subsection*{Expected Deliverables}
\begin{itemize}[leftmargin=*,itemsep=0.2em]
\item State one theorem-shaped candidate centered on 'barrier-aware transfer' satisfying 'Either construct a local reduction family G satisfying locality bounds L and preserving witness structure W, or prove that any such G would contradict barrier theorem B.'.
\item Make 'A barrier-aware reduction gadget family or an impossibility theorem explaining why the reduction regime cannot be crossed.' explicit enough that an expert could inspect the object or transfer law.
\item Record one decisive negative route explaining what would fail if this discovery layer does not land.
\end{itemize}
\subsection*{Stop Criteria}
\begin{itemize}[leftmargin=*,itemsep=0.2em]
\item Stop immediately if the branch collapses into the anti-pattern 'ignoring the known local-reduction barrier theorem'.
\item Stop if the branch yields no theorem-shaped claim, carrier, transport principle, invariant, or representation change.
\item Stop after one explicit positive candidate and one explicit failure probe have both been recorded.
\item Escalate if 'barrier-aware transfer' remains purely schematic after the first explicit object or theorem candidate.
\end{itemize}
\subsection*{External Framework Search Heuristic}
\textbf{Recommended.} True

\textbf{Mode.} targeted-foundation-search

\textbf{Justification.} Search externally for a mathematical framework with similar structural properties only after local routes, InformalMathLib, and explicit negative probes fail to provide the missing carrier, invariant, transfer law, or correspondence theorem.

\textbf{Triggers.} cross-family-pressure, bridge-family-foundation-gap

\textbf{Candidate framework regions.} lower-bound barrier frameworks, local-reduction impossibility theory, gadget and locality-preservation frameworks

\subsection*{Local Prerequisites Before Browsing}
\begin{itemize}[leftmargin=*,itemsep=0.2em]
\item Record the failed local theorem lookup, InformalMathLib lookup, and negative probe that justify leaving local search.
\item State the missing operation concretely: carrier realization, transport, alignment, explicitness, characterization, or invariant discovery.
\item Reject at least one purely topical analogy before searching for an external framework.
\end{itemize}
\subsection*{Query Templates}
\begin{itemize}[leftmargin=*,itemsep=0.2em]
\item mathematical framework for Barrier-Aware Transfer Prompt with preserved invariant and bridge theorem
\item Barrier-Aware Transfer Prompt correspondence theorem transport invariant
\item Barrier-Aware Transfer Prompt framework Design a transfer mechanism that survives known local-reduction barriers, or prove that no such mechanism exists under explicit assumptions. theorem
\item theorem or framework matching: Either construct a local reduction family G satisfying locality bounds L and preserving witness structure W, or prove that any such G would contradict barrier theorem B.
\end{itemize}
\subsection*{Required Questions}
\begin{itemize}[leftmargin=*,itemsep=0.2em]
\item What exact barrier theorem constrains the route?
\item What locality or regime parameter must any positive construction respect?
\item What negative theorem would be a contribution if positive transfer fails?
\end{itemize}
\subsection*{Search Modes}
\begin{itemize}[leftmargin=*,itemsep=0.2em]
\item barrier-aware gadget design
\item impossibility theorem search
\item reduction regime decomposition
\end{itemize}
\section{disclayer.p-vs-np.pattern.problem-representation-shift}
\textbf{Kind.} epistemic-pattern

\textbf{Target.} problem-representation-shift

\textbf{Routing priority.} 6.548

\textbf{Dispatch rank.} 11

\textbf{Spawn reason.} Spawn 'problem-representation-shift' as a cross-cutting discovery layer because the blocked frontier spans multiple breakthrough families and needs a broader reframing.

\textbf{Rationale.} Several blocked routes look misrepresented rather than simply false, so representation change is itself a discovery pattern.

\subsection*{System Prompt}
You are testing whether the blocked route is misrepresented rather than false. Search for an alternative encoding, normal form, dual object, or reformulation in which the route becomes mathematically transparent.

\subsection*{User Task}
Propose representation changes that preserve the source-side content but expose a better target interface, and state the first theorem each representation would make possible.

\subsection*{Expected Deliverables}
\begin{itemize}[leftmargin=*,itemsep=0.2em]
\item Produce one reframing under 'problem-representation-shift' that changes the search state rather than restating the blocked lemma.
\item Connect the reframing back to a theorem-shaped objective compatible with 'There exists an admissible MCSP witness family, realized as a family of instances equipped with a monotone witness measure, such that whenever A barely-nontrivial lower bound for an explicit problem Q in a weak model holds for MCSP in nonuniform-circuits, the object organizes the source-side mathematical content while preserving the invariant MCSP complexity profile in nonuniform-circuits and yields major consequences such as NP not in P/poly.'.
\item Record one stop condition showing when the pattern has become superficial rather than discovery-bearing.
\end{itemize}
\subsection*{Stop Criteria}
\begin{itemize}[leftmargin=*,itemsep=0.2em]
\item Stop immediately if the branch collapses into the anti-pattern 'adding a stronger carrier without specifying the target object class map'.
\item Stop if the branch yields no theorem-shaped claim, carrier, transport principle, invariant, or representation change.
\item Stop after one explicit positive candidate and one explicit failure probe have both been recorded.
\item Escalate if 'problem-representation-shift' only re-labels the blocked frontier instead of reducing proof debt or changing the representation.
\end{itemize}
\subsection*{External Framework Search Heuristic}
\textbf{Recommended.} True

\textbf{Mode.} targeted-foundation-search

\textbf{Justification.} Search externally for a mathematical framework with similar structural properties only after local routes, InformalMathLib, and explicit negative probes fail to provide the missing carrier, invariant, transfer law, or correspondence theorem.

\textbf{Triggers.} structural-analogy, cross-family-pressure, pattern-level-reframing, bridge-family-foundation-gap

\textbf{Candidate framework regions.} normal-form frameworks, equivalent reformulation theorems, model/encoding correspondence theories, proof complexity, pseudorandomness / derandomization

\subsection*{Local Prerequisites Before Browsing}
\begin{itemize}[leftmargin=*,itemsep=0.2em]
\item Record the failed local theorem lookup, InformalMathLib lookup, and negative probe that justify leaving local search.
\item State the missing operation concretely: carrier realization, transport, alignment, explicitness, characterization, or invariant discovery.
\item Reject at least one purely topical analogy before searching for an external framework.
\end{itemize}
\subsection*{Query Templates}
\begin{itemize}[leftmargin=*,itemsep=0.2em]
\item mathematical framework for Problem Representation Shift Prompt with preserved invariant and bridge theorem
\item Problem Representation Shift Prompt correspondence theorem transport invariant
\item Problem Representation Shift Prompt framework Several blocked routes look misrepresented rather than simply false, so representation change is itself a discovery pattern. theorem
\item theorem or framework matching: There exists an admissible MCSP witness family, realized as a family of instances equipped with a monotone witness measure, such that whenever A barely-nontrivial lower bound for an explicit problem Q in a weak model holds for MCSP in nonuniform-circuits, the object organizes the source-side mathematical content while preserving the invariant MCSP complexity profile in nonuniform-circuits and yields major consequences such as NP not in P/poly.
\end{itemize}
\subsection*{Required Questions}
\begin{itemize}[leftmargin=*,itemsep=0.2em]
\item Which part of the current representation hides the real inferential structure?
\item What alternative object language or normal form would expose the missing bridge?
\item What mathematical content must stay invariant under the representation change?
\end{itemize}
\subsection*{Search Modes}
\begin{itemize}[leftmargin=*,itemsep=0.2em]
\item bridge split
\item counterfactual sharpening
\item cross-model transport principle
\item dual characterization of the target object
\item duality or normal form search
\item generalization by faithful embedding
\item invariant-preserving compilation theorem
\item normal-form mediation
\item object-class correspondence search
\item parameter-preserving reformulation
\item problem-specific to generic theorem bridge
\item representation normal form for formula-size transfer
\item semantic correspondence theorem
\item two-way implication decomposition
\item typed abstraction search
\end{itemize}
\section{disclayer.p-vs-np.family.alignment}
\textbf{Kind.} breakthrough-family

\textbf{Target.} alignment

\textbf{Routing priority.} 2.222

\textbf{Dispatch rank.} 12

\textbf{Spawn reason.} Spawn a dedicated 'alignment' discovery layer because obstruction 'target-alignment' carries repeated pressure across the blocked frontier.

\textbf{Rationale.} Identify or prove an equivalence between the current carrier class and the object class required by the target theorem.

\subsection*{System Prompt}
You are aligning a local carrier class with the object class demanded by an existing target theorem. Topical similarity is not enough. Search for equivalence, correspondence, or a mediating normal form.

\subsection*{User Task}
Describe the object-class map from the current carrier to the target theorem's object class, including the quantitative property that must survive the correspondence.

\subsection*{Expected Deliverables}
\begin{itemize}[leftmargin=*,itemsep=0.2em]
\item State one theorem-shaped candidate centered on 'alignment' satisfying 'Every admissible carrier C from the local route induces an object D in the target theorem's class, and D preserves the quantitative property needed to activate the target theorem.'.
\item Make 'An alignment theorem, correspondence theorem, or mediating normal form linking the local carrier to the target's sparse-problem or distinguisher objects.' explicit enough that an expert could inspect the object or transfer law.
\item Record one decisive negative route explaining what would fail if this discovery layer does not land.
\end{itemize}
\subsection*{Stop Criteria}
\begin{itemize}[leftmargin=*,itemsep=0.2em]
\item Stop immediately if the branch collapses into the anti-pattern 'treating topical similarity as object equivalence'.
\item Stop if the branch yields no theorem-shaped claim, carrier, transport principle, invariant, or representation change.
\item Stop after one explicit positive candidate and one explicit failure probe have both been recorded.
\item Escalate if 'alignment' remains purely schematic after the first explicit object or theorem candidate.
\end{itemize}
\subsection*{External Framework Search Heuristic}
\textbf{Recommended.} True

\textbf{Mode.} targeted-foundation-search

\textbf{Justification.} Search externally for a mathematical framework with similar structural properties only after local routes, InformalMathLib, and explicit negative probes fail to provide the missing carrier, invariant, transfer law, or correspondence theorem.

\textbf{Triggers.} structural-analogy, cross-family-pressure, bridge-family-foundation-gap

\textbf{Candidate framework regions.} sparse-object and distinguisher frameworks, duality or correspondence theorems, representation-theoretic normal forms

\subsection*{Local Prerequisites Before Browsing}
\begin{itemize}[leftmargin=*,itemsep=0.2em]
\item Record the failed local theorem lookup, InformalMathLib lookup, and negative probe that justify leaving local search.
\item State the missing operation concretely: carrier realization, transport, alignment, explicitness, characterization, or invariant discovery.
\item Reject at least one purely topical analogy before searching for an external framework.
\end{itemize}
\subsection*{Query Templates}
\begin{itemize}[leftmargin=*,itemsep=0.2em]
\item mathematical framework for Alignment Theorem Prompt with preserved invariant and bridge theorem
\item Alignment Theorem Prompt correspondence theorem transport invariant
\item Alignment Theorem Prompt framework Identify or prove an equivalence between the current carrier class and the object class required by the target theorem. theorem
\item theorem or framework matching: Every admissible carrier C from the local route induces an object D in the target theorem's class, and D preserves the quantitative property needed to activate the target theorem.
\end{itemize}
\subsection*{Required Questions}
\begin{itemize}[leftmargin=*,itemsep=0.2em]
\item What is the target theorem's exact object class?
\item How does the local carrier induce an object in that class?
\item Which parameter or sparsity threshold is hardest to preserve?
\end{itemize}
\subsection*{Search Modes}
\begin{itemize}[leftmargin=*,itemsep=0.2em]
\item object-class correspondence search
\item normal-form mediation
\item dual characterization of the target object
\end{itemize}
\section{disclayer.p-vs-np.family.explicitness}
\textbf{Kind.} breakthrough-family

\textbf{Target.} explicitness

\textbf{Routing priority.} 2.107

\textbf{Dispatch rank.} 13

\textbf{Spawn reason.} Spawn a dedicated 'explicitness' discovery layer because obstruction 'pit-upgrade' carries repeated pressure across the blocked frontier.

\textbf{Rationale.} Upgrade a semantic or existential algebraic object into an explicit algorithmic derandomization premise.

\subsection*{System Prompt}
You are upgrading a semantic or existential object into an explicit algorithmic one. Do not accept existence without uniformity, efficiency, or extractability.

\subsection*{User Task}
Search for an explicit generator, certificate family, extractor, or uniform construction rule that turns the blocked semantic object into a derandomization-ready or algorithmically inspectable premise.

\subsection*{Expected Deliverables}
\begin{itemize}[leftmargin=*,itemsep=0.2em]
\item State one theorem-shaped candidate centered on 'explicitness' satisfying 'If semantic object C exists for the relevant circuit class, then one can uniformly and efficiently extract an explicit structure G that yields deterministic PIT for that class.'.
\item Make 'An explicit generator, certificate family, or algebraic witness system that is constructive enough for PIT rather than merely semantic.' explicit enough that an expert could inspect the object or transfer law.
\item Record one decisive negative route explaining what would fail if this discovery layer does not land.
\end{itemize}
\subsection*{Stop Criteria}
\begin{itemize}[leftmargin=*,itemsep=0.2em]
\item Stop immediately if the branch collapses into the anti-pattern 'equating a black-box hitting-set statement with an explicit generator'.
\item Stop if the branch yields no theorem-shaped claim, carrier, transport principle, invariant, or representation change.
\item Stop after one explicit positive candidate and one explicit failure probe have both been recorded.
\item Escalate if 'explicitness' remains purely schematic after the first explicit object or theorem candidate.
\end{itemize}
\subsection*{External Framework Search Heuristic}
\textbf{Recommended.} True

\textbf{Mode.} targeted-foundation-search

\textbf{Justification.} Search externally for a mathematical framework with similar structural properties only after local routes, InformalMathLib, and explicit negative probes fail to provide the missing carrier, invariant, transfer law, or correspondence theorem.

\textbf{Triggers.} structural-analogy, cross-family-pressure, bridge-family-foundation-gap

\textbf{Candidate framework regions.} derandomization and generator frameworks, algebraic certificate constructions, extractor / hitter frameworks

\subsection*{Local Prerequisites Before Browsing}
\begin{itemize}[leftmargin=*,itemsep=0.2em]
\item Record the failed local theorem lookup, InformalMathLib lookup, and negative probe that justify leaving local search.
\item State the missing operation concretely: carrier realization, transport, alignment, explicitness, characterization, or invariant discovery.
\item Reject at least one purely topical analogy before searching for an external framework.
\end{itemize}
\subsection*{Query Templates}
\begin{itemize}[leftmargin=*,itemsep=0.2em]
\item mathematical framework for Explicitness Upgrade Prompt with preserved invariant and bridge theorem
\item Explicitness Upgrade Prompt correspondence theorem transport invariant
\item Explicitness Upgrade Prompt framework Upgrade a semantic or existential algebraic object into an explicit algorithmic derandomization premise. theorem
\item theorem or framework matching: If semantic object C exists for the relevant circuit class, then one can uniformly and efficiently extract an explicit structure G that yields deterministic PIT for that class.
\end{itemize}
\subsection*{Required Questions}
\begin{itemize}[leftmargin=*,itemsep=0.2em]
\item What part of the current object is only existential or semantic?
\item What explicit structure would replace that existential step?
\item What efficiency class is required for the upgraded object to matter?
\end{itemize}
\subsection*{Search Modes}
\begin{itemize}[leftmargin=*,itemsep=0.2em]
\item explicit extractor or generator search
\item certificate map construction
\item semantic-to-algorithmic upgrade theorem
\end{itemize}
\section{disclayer.p-vs-np.family.constructive-instantiation}
\textbf{Kind.} breakthrough-family

\textbf{Target.} constructive instantiation

\textbf{Routing priority.} 1.966

\textbf{Dispatch rank.} 14

\textbf{Spawn reason.} Spawn a dedicated 'constructive instantiation' discovery layer because obstruction 'schema-to-examples' carries repeated pressure across the blocked frontier.

\textbf{Rationale.} Turn an abstract schema into explicit natural example families or prove that the schema cannot instantiate naturally in the current way.

\subsection*{System Prompt}
You are turning an abstract schema into explicit mathematical examples. Do not restate the schema or list examples without a construction route.

\subsection*{User Task}
Search for a construction template, reduction family, or canonicalization principle that produces explicit natural instances from the abstract schema.

\subsection*{Expected Deliverables}
\begin{itemize}[leftmargin=*,itemsep=0.2em]
\item State one theorem-shaped candidate centered on 'constructive instantiation' satisfying 'From abstract schema S, construct an explicit family of natural instances E\_n satisfying properties P and realizing target theorem T.'.
\item Make 'A construction template, example-generating reduction family, or a canonicalization principle for natural examples.' explicit enough that an expert could inspect the object or transfer law.
\item Record one decisive negative route explaining what would fail if this discovery layer does not land.
\end{itemize}
\subsection*{Stop Criteria}
\begin{itemize}[leftmargin=*,itemsep=0.2em]
\item Stop immediately if the branch collapses into the anti-pattern 'restating the schema as if it already supplied examples'.
\item Stop if the branch yields no theorem-shaped claim, carrier, transport principle, invariant, or representation change.
\item Stop after one explicit positive candidate and one explicit failure probe have both been recorded.
\item Escalate if 'constructive instantiation' remains purely schematic after the first explicit object or theorem candidate.
\end{itemize}
\subsection*{External Framework Search Heuristic}
\textbf{Recommended.} True

\textbf{Mode.} targeted-foundation-search

\textbf{Justification.} Search externally for a mathematical framework with similar structural properties only after local routes, InformalMathLib, and explicit negative probes fail to provide the missing carrier, invariant, transfer law, or correspondence theorem.

\textbf{Triggers.} cross-family-pressure

\textbf{Candidate framework regions.} canonical example construction frameworks, natural proof or canonicalization methods, explicit reduction template families

\subsection*{Local Prerequisites Before Browsing}
\begin{itemize}[leftmargin=*,itemsep=0.2em]
\item Record the failed local theorem lookup, InformalMathLib lookup, and negative probe that justify leaving local search.
\item State the missing operation concretely: carrier realization, transport, alignment, explicitness, characterization, or invariant discovery.
\item Reject at least one purely topical analogy before searching for an external framework.
\end{itemize}
\subsection*{Query Templates}
\begin{itemize}[leftmargin=*,itemsep=0.2em]
\item mathematical framework for Constructive Instantiation Prompt with preserved invariant and bridge theorem
\item Constructive Instantiation Prompt correspondence theorem transport invariant
\item Constructive Instantiation Prompt framework Turn an abstract schema into explicit natural example families or prove that the schema cannot instantiate naturally in the current way. theorem
\item theorem or framework matching: From abstract schema S, construct an explicit family of natural instances E\_n satisfying properties P and realizing target theorem T.
\end{itemize}
\subsection*{Required Questions}
\begin{itemize}[leftmargin=*,itemsep=0.2em]
\item What is the first explicit instance family that should be produced?
\item What reduction or construction template yields that family?
\item If the construction fails, what does that say about the schema's naturality?
\end{itemize}
\subsection*{Search Modes}
\begin{itemize}[leftmargin=*,itemsep=0.2em]
\item canonical example construction
\item template instantiation
\item example-preserving reduction search
\end{itemize}
\section{disclayer.p-vs-np.family.semantic-bridge}
\textbf{Kind.} breakthrough-family

\textbf{Target.} semantic bridge

\textbf{Routing priority.} 1.946

\textbf{Dispatch rank.} 15

\textbf{Spawn reason.} Spawn a dedicated 'semantic bridge' discovery layer because obstruction 'characterization-bridge' carries repeated pressure across the blocked frontier.

\textbf{Rationale.} Bridge a local definitional object to a global characterization theorem.

\subsection*{System Prompt}
You are bridging a local definitional object to a global characterization theorem. Syntactic normalization is insufficient. Search for the missing semantic equivalence, duality, or witness-extraction principle.

\subsection*{User Task}
Decompose the biconditional or characterization gap into the minimal semantic implications that would connect the local definition to the global theorem.

\subsection*{Expected Deliverables}
\begin{itemize}[leftmargin=*,itemsep=0.2em]
\item State one theorem-shaped candidate centered on 'semantic bridge' satisfying 'Local definitional object D has property P if and only if global characterization condition Q holds, with both directions proved at the right level of abstraction.'.
\item Make 'A semantic correspondence, duality, witness extraction principle, or completeness notion that turns the local definition into the target biconditional.' explicit enough that an expert could inspect the object or transfer law.
\item Record one decisive negative route explaining what would fail if this discovery layer does not land.
\end{itemize}
\subsection*{Stop Criteria}
\begin{itemize}[leftmargin=*,itemsep=0.2em]
\item Stop immediately if the branch collapses into the anti-pattern 'normalizing syntax without adding semantic force'.
\item Stop if the branch yields no theorem-shaped claim, carrier, transport principle, invariant, or representation change.
\item Stop after one explicit positive candidate and one explicit failure probe have both been recorded.
\item Escalate if 'semantic bridge' remains purely schematic after the first explicit object or theorem candidate.
\end{itemize}
\subsection*{External Framework Search Heuristic}
\textbf{Recommended.} True

\textbf{Mode.} targeted-foundation-search

\textbf{Justification.} Search externally for a mathematical framework with similar structural properties only after local routes, InformalMathLib, and explicit negative probes fail to provide the missing carrier, invariant, transfer law, or correspondence theorem.

\textbf{Triggers.} structural-analogy, cross-family-pressure, bridge-family-foundation-gap

\textbf{Candidate framework regions.} duality frameworks, characterization theorems in proof complexity, semantic completeness correspondences

\subsection*{Local Prerequisites Before Browsing}
\begin{itemize}[leftmargin=*,itemsep=0.2em]
\item Record the failed local theorem lookup, InformalMathLib lookup, and negative probe that justify leaving local search.
\item State the missing operation concretely: carrier realization, transport, alignment, explicitness, characterization, or invariant discovery.
\item Reject at least one purely topical analogy before searching for an external framework.
\end{itemize}
\subsection*{Query Templates}
\begin{itemize}[leftmargin=*,itemsep=0.2em]
\item mathematical framework for Semantic Bridge Prompt with preserved invariant and bridge theorem
\item Semantic Bridge Prompt correspondence theorem transport invariant
\item Semantic Bridge Prompt framework Bridge a local definitional object to a global characterization theorem. theorem
\item theorem or framework matching: Local definitional object D has property P if and only if global characterization condition Q holds, with both directions proved at the right level of abstraction.
\end{itemize}
\subsection*{Required Questions}
\begin{itemize}[leftmargin=*,itemsep=0.2em]
\item What semantic content is missing from the current definition-level object?
\item Which direction of the characterization is harder and why?
\item What duality or witness principle could mediate the bridge?
\end{itemize}
\subsection*{Search Modes}
\begin{itemize}[leftmargin=*,itemsep=0.2em]
\item semantic correspondence theorem
\item duality or normal form search
\item two-way implication decomposition
\end{itemize}
\section{disclayer.p-vs-np.family.categorical-reframing}
\textbf{Kind.} breakthrough-family

\textbf{Target.} categorical-reframing

\textbf{Routing priority.} 0.910

\textbf{Dispatch rank.} 16

\textbf{Spawn reason.} Spawn a dedicated 'categorical-reframing' discovery layer because obstruction 'structural-composition-gap' carries repeated pressure across the blocked frontier.

\textbf{Rationale.} Recast the blocked frontier as a diagrammatic composition problem: identify objects, morphisms, mediating objects, and the exact commutativity or factorization obligation.

\subsection*{System Prompt}
Use category theory as a structural and compositional discovery lens, not as decorative abstraction. Identify the source object, target object, typed morphisms already available, and the diagram that should commute if the route is real. Search for a mediating object with a universal property, test whether the route factors through it, and classify failure as no morphism, non-functorial transport, broken naturality, missing factorization, or failed universal property.

\subsection*{User Task}
Reframe the blocked frontier as a family of diagrammatic composition problems. State candidate objects, morphisms, commutative squares, factorization claims, and universal properties that would make the live routes structurally compose.

\subsection*{Expected Deliverables}
\begin{itemize}[leftmargin=*,itemsep=0.2em]
\item State one theorem-shaped candidate centered on 'categorical-reframing' satisfying 'There exists a mediating object M and typed morphisms making the route commute; otherwise the failure is certified as no morphism, non-functorial transport, broken naturality, or missing universal property.'.
\item Make 'A structural account of the route in terms of source and target objects, typed morphisms, a mediating universal object, and the factorization or naturality law that would make the route compose.' explicit enough that an expert could inspect the object or transfer law.
\item Record one decisive negative route explaining what would fail if this discovery layer does not land.
\end{itemize}
\subsection*{Stop Criteria}
\begin{itemize}[leftmargin=*,itemsep=0.2em]
\item Stop immediately if the branch collapses into the anti-pattern 'renaming route components as categorical objects without reducing proof debt'.
\item Stop if the branch yields no theorem-shaped claim, carrier, transport principle, invariant, or representation change.
\item Stop after one explicit positive candidate and one explicit failure probe have both been recorded.
\item Escalate if 'categorical-reframing' remains purely schematic after the first explicit object or theorem candidate.
\end{itemize}
\subsection*{External Framework Search Heuristic}
\textbf{Recommended.} True

\textbf{Mode.} targeted-foundation-search

\textbf{Justification.} Search externally for a mathematical framework with similar structural properties only after local routes, InformalMathLib, and explicit negative probes fail to provide the missing carrier, invariant, transfer law, or correspondence theorem.

\textbf{Triggers.} cross-family-pressure, bridge-family-foundation-gap

\textbf{Candidate framework regions.} sparse-object and distinguisher frameworks, duality or correspondence theorems, representation-theoretic normal forms, lower-bound barrier frameworks, local-reduction impossibility theory

\subsection*{Local Prerequisites Before Browsing}
\begin{itemize}[leftmargin=*,itemsep=0.2em]
\item Record the failed local theorem lookup, InformalMathLib lookup, and negative probe that justify leaving local search.
\item State the missing operation concretely: carrier realization, transport, alignment, explicitness, characterization, or invariant discovery.
\item Reject at least one purely topical analogy before searching for an external framework.
\end{itemize}
\subsection*{Query Templates}
\begin{itemize}[leftmargin=*,itemsep=0.2em]
\item mathematical framework for Categorical Reframing Prompt with preserved invariant and bridge theorem
\item Categorical Reframing Prompt correspondence theorem transport invariant
\item Categorical Reframing Prompt framework Recast the blocked frontier as a diagrammatic composition problem: identify objects, morphisms, mediating objects, and the exact commutativity or factorization obligation. theorem
\item theorem or framework matching: There exists a mediating object M and typed morphisms making the route commute; otherwise the failure is certified as no morphism, non-functorial transport, broken naturality, or missing universal property.
\end{itemize}
\subsection*{Required Questions}
\begin{itemize}[leftmargin=*,itemsep=0.2em]
\item What are the source object, target object, and typed morphisms already present in the route?
\item What diagram should commute if the route is legitimate?
\item What mediating object would be universal for this route?
\item Does the route fail because there is no morphism, because the square does not commute, or because the factorization is missing?
\end{itemize}
\subsection*{Search Modes}
\begin{itemize}[leftmargin=*,itemsep=0.2em]
\item diagram construction
\item commutativity checking
\item factorization search
\item universal-property search
\item adjunction-style dualization
\end{itemize}
\section{disclayer.p-vs-np.family.equivalence-coherence-reframing}
\textbf{Kind.} breakthrough-family

\textbf{Target.} equivalence-coherence-reframing

\textbf{Routing priority.} 0.880

\textbf{Dispatch rank.} 17

\textbf{Spawn reason.} Spawn a dedicated 'equivalence-coherence-reframing' discovery layer because obstruction 'equivalence-coherence-gap' carries repeated pressure across the blocked frontier.

\textbf{Rationale.} Recast the blocked frontier as an identity, equivalence, transport, or coherence problem rather than as a missing brute-force bridge.

\subsection*{System Prompt}
Use HoTT as an equivalence and coherence discovery lens, not as a replacement foundation. Type the active concepts and bridge objects, ask whether two formulations are equal, equivalent, or only analogous, search for transportable invariants, and classify failure as no equivalence witness, incoherent transport, or weak equivalence where strict identity is needed.

\subsection*{User Task}
Reframe the blocked frontier as a problem about sameness, equivalence, transport, and coherence. State the candidate equivalence claims, the transport obligations they would induce, and the coherence failures that currently block the route.

\subsection*{Expected Deliverables}
\begin{itemize}[leftmargin=*,itemsep=0.2em]
\item State one theorem-shaped candidate centered on 'equivalence-coherence-reframing' satisfying 'Either exhibit an equivalence E between the live formulations and prove that the target-relevant invariant transports coherently along E, or prove that no such equivalence witness exists in the current regime.'.
\item Make 'A typed account of active formulations as types and terms, an explicit equivalence witness or non-equivalence diagnosis, and the transport law needed to move invariants and theorem hypotheses across that equivalence.' explicit enough that an expert could inspect the object or transfer law.
\item Record one decisive negative route explaining what would fail if this discovery layer does not land.
\end{itemize}
\subsection*{Stop Criteria}
\begin{itemize}[leftmargin=*,itemsep=0.2em]
\item Stop immediately if the branch collapses into the anti-pattern 'treating topical similarity as equivalence'.
\item Stop if the branch yields no theorem-shaped claim, carrier, transport principle, invariant, or representation change.
\item Stop after one explicit positive candidate and one explicit failure probe have both been recorded.
\item Escalate if 'equivalence-coherence-reframing' remains purely schematic after the first explicit object or theorem candidate.
\end{itemize}
\subsection*{External Framework Search Heuristic}
\textbf{Recommended.} True

\textbf{Mode.} targeted-foundation-search

\textbf{Justification.} Search externally for a mathematical framework with similar structural properties only after local routes, InformalMathLib, and explicit negative probes fail to provide the missing carrier, invariant, transfer law, or correspondence theorem.

\textbf{Triggers.} structural-analogy, cross-family-pressure, bridge-family-foundation-gap

\textbf{Candidate framework regions.} sparse-object and distinguisher frameworks, duality or correspondence theorems, representation-theoretic normal forms, derandomization and generator frameworks, algebraic certificate constructions

\subsection*{Local Prerequisites Before Browsing}
\begin{itemize}[leftmargin=*,itemsep=0.2em]
\item Record the failed local theorem lookup, InformalMathLib lookup, and negative probe that justify leaving local search.
\item State the missing operation concretely: carrier realization, transport, alignment, explicitness, characterization, or invariant discovery.
\item Reject at least one purely topical analogy before searching for an external framework.
\end{itemize}
\subsection*{Query Templates}
\begin{itemize}[leftmargin=*,itemsep=0.2em]
\item mathematical framework for Equivalence / Coherence Reframing Prompt with preserved invariant and bridge theorem
\item Equivalence / Coherence Reframing Prompt correspondence theorem transport invariant
\item Equivalence / Coherence Reframing Prompt framework Recast the blocked frontier as an identity, equivalence, transport, or coherence problem rather than as a missing brute-force bridge. theorem
\item theorem or framework matching: Either exhibit an equivalence E between the live formulations and prove that the target-relevant invariant transports coherently along E, or prove that no such equivalence witness exists in the current regime.
\end{itemize}
\subsection*{Required Questions}
\begin{itemize}[leftmargin=*,itemsep=0.2em]
\item Which active concepts should be treated as types, and which concrete witnesses should be treated as terms?
\item Are the live operationalizations equal, equivalent, incomparable, or only analogous?
\item What invariant or theorem hypothesis must transport along the proposed equivalence?
\item Where does coherence fail when the same route is identified in two different ways?
\end{itemize}
\subsection*{Search Modes}
\begin{itemize}[leftmargin=*,itemsep=0.2em]
\item equivalence witness search
\item transport-law search
\item coherence checking
\item identity-vs-equivalence discrimination
\end{itemize}

\end{document}
